% Generated draft: requires mathematical and source review.
\documentclass[11pt]{article}
\usepackage[margin=25mm]{geometry}
\usepackage{fontspec}
\setmainfont{Latin Modern Roman}
\usepackage{amsmath,amssymb,mathtools,unicode-math}
\setmathfont{Latin Modern Math}
\usepackage{enumitem,needspace,xurl,hyperref}
\setlength{\parindent}{0pt}
\setlength{\parskip}{5pt}
\setlength{\emergencystretch}{3em}
\newcommand{\N}{\mathbb{N}}
\newcommand{\Z}{\mathbb{Z}}
\newcommand{\Q}{\mathbb{Q}}
\newcommand{\R}{\mathbb{R}}
\newcommand{\C}{\mathbb{C}}
\newcommand{\F}{\mathbb{F}}
\newcommand{\A}{\mathbb{A}}
\newcommand{\PP}{\mathbb{P}}
\newcommand{\E}{\mathbb{E}}
\newcommand{\eps}{\varepsilon}
\newcommand{\dd}{\,\mathrm{d}}
\begin{document}
\section*{UnsolvedMath ID: 30004992}
\textbf{Source problem code:} catalogue entry
\subsection{Definitions and mathematical statement}
Let \(G\) be a second-countable locally compact Hausdorff topological group. A unitary representation is understood to be strongly continuous. The group \(G\) is \emph{type I} if, for every unitary representation \(\pi\) for which the von Neumann algebra \(\pi(G)''\) is a factor, that factor is of type I, that is, isomorphic to \(B(\mathcal K)\) for some Hilbert space \(\mathcal K\).

A locally compact group \(H\) is \emph{amenable} if there is a positive linear functional \(m:L^\infty(H)\to\C\), with respect to a left Haar measure, such that \(m(1)=1\) and \(m(f\circ L_h)=m(f)\) for all \(h\in H\) and \(f\in L^\infty(H)\), where \(L_h(x)=hx\). A closed subgroup \(H\leq G\) is \emph{cocompact} if the quotient space \(G/H\) is compact. We use the usual interpretation that the sought subgroup is closed and carries its subspace topology; the supplied statement does not explicitly include the word ``closed.''

\textbf{Main conjecture (Conjecture 1 in the supplied excerpt).} For every second-countable locally compact group \(G\),
\[
 G\text{ is type I}
 \quad\Longrightarrow\quad
 \exists\,H\leq G\;\bigl(H\text{ is closed and amenable, and }G/H\text{ is compact}\bigr).
\]
No nonamenability assumption is imposed. If \(G\) is amenable, the conclusion holds by taking \(H=G\); this explains the source's description of the statement as ``meaningless'' in that case.

\textbf{Qualification concerning a converse.} The excerpt observes that amenable groups need not be type I, so the unrestricted converse fails. It only suggests that a converse might hold under the additional assumption that \(G\) has trivial amenable radical, meaning that its only amenable closed normal subgroup is \(\{e\}\). In precise conditional form, the suggested question is whether, for a second-countable locally compact \(G\) with this property, the existence of a closed cocompact amenable subgroup implies that \(G\) is type I. This suggestion is separate from the main conjecture.

\textbf{Additional tree context in the supplied excerpt.} Let \(T\) be a locally finite tree in which every vertex has degree at least \(3\). Its end space \(\partial T\), consisting of equivalence classes of infinite rays that eventually have a common tail, has its usual compact Hausdorff topology. Let \(G\) be a noncompact closed subgroup of \(\operatorname{Aut}(T)\), with the topology of pointwise convergence on vertices, and suppose that its action on \(\partial T\) is minimal, meaning that every orbit is dense. A group is \emph{liminal} if every irreducible unitary representation \(\pi\) maps its full group \(C^*\)-algebra onto the compact operators on the representation space. The separately stated Conjecture 2, attributed in the excerpt to Nebbia [8] and Houdayer--Raum [7], asserts the equivalence of
\[
\begin{aligned}
\text{(i)}\;&G\text{ is liminal},\\
\text{(ii)}\;&G\text{ is type I},\\
\text{(iii)}\;&G\text{ acts transitively on ordered pairs of distinct points of }\partial T.
\end{aligned}
\]
The excerpt reports that, under these tree hypotheses, (iii) is equivalent to the existence of a cocompact amenable subgroup, citing its reference [3]. It states that (i) implies (ii) in general, attributes the conjectural implication (iii) to (i) to Nebbia [8], and attributes the conjectural implication (ii) to (iii) to Houdayer--Raum [7]. These are contextual source assertions, not additional hypotheses of the main conjecture.

The excerpt also reports the following theorem of Houdayer--Raum [7]: for \(G\) and \(T\) as above, if the group von Neumann algebra \(L(G)=\lambda_G(G)''\), generated by the left regular representation on \(L^2(G)\), is amenable as a von Neumann algebra, then every vertex stabilizer \(G_v\) acts \(2\)-transitively on the neighbors of \(v\). It gives the type I property of \(L(G)\) as a sufficient condition for this operator-algebraic amenability. The excerpt further states that boundary \(2\)-transitivity implies this local \(2\)-transitivity and that the converse need not hold. The numbered references underlying these contextual claims are not identified bibliographically in the supplied record. A subsequent mention of a conjectural description of nondiscrete type I groups in reference [4] supplies no formulation and therefore cannot be reconstructed here.

\textbf{Supplied status qualification.} The registry labels the problem partially solved. Its literature assessment dated 2026-08-22 says that the conjecture is proved in important hyperbolic and unimodular cases and remains open in general. The precise hypotheses of those cases are not supplied. This is dated supplied curation, not an independently verified mathematical or literature assessment.
\subsection{Short English statement}
The main question asks whether every second-countable locally compact type I group contains a closed amenable subgroup with compact quotient. For amenable groups, choosing the whole group already gives such a subgroup. The supplied excerpt also discusses a possible converse under a trivial amenable radical and a separate conjecture about groups acting on trees. Its partially solved label and claims about known cases are supplied dated assessments requiring review.
\subsection{Sources}
\begin{itemize}
\item Local registry: \texttt{data/unsolvedmath/problems.json}, supplied record with permanent UnsolvedMath ID \(30004992\), source problem code \texttt{catalogue entry}, title \emph{Cocompact Amenable Subgroups of Type-I Groups}, in the Oberwolfach Reports Open Problems collection. No proposer is named; the supplied proposed year is 2022.
\item Supplied report citation: \emph{Groups and Dynamics: Topology, Measure, and Borel Structure} (2023), Oberwolfach Reports, DOI \texttt{10.4171/owr/2022/3}, pp.~107--169. \url{https://doi.org/10.4171/owr/2022/3}. The pasted excerpt includes a page marker for p.~125. The record labels the statement \texttt{corrected\_verified} and says it was reconstructed and checked against the official report; that verification is supplied provenance and has not been independently repeated here.
\item Supplied literature reference: Caprace--Kalantar--Monod, \emph{A type I conjecture and boundary representations}. \url{https://arxiv.org/abs/2110.00190}.
\item Supplied literature reference: Glasner, \emph{Boundary representations of locally compact hyperbolic groups}. \url{https://doi.org/10.1112/plms.70100}.
\item Supplied literature check date: \texttt{2026-08-22}. The assessment describes the problem as partially solved, with important hyperbolic and unimodular cases proved and the general conjecture open. This is supplied dated curation derived from the OWR record, not independently verified; the record expressly requests independent verification of time-sensitive status and cited claims.
\item Contextual attribution in the excerpt: Nebbia [8], Houdayer--Raum [7], and unidentified references [3] and [4]. The supplied record does not provide the corresponding bibliography, and no identification of these numbered references with the literature links above is assumed.
\end{itemize}
\end{document}
